% =====================================================================
% 04_prop1.tex — Proposition 1: prices and production (WP-B)
% =====================================================================
\section{Proposition 1: effects of a safety shock on prices and production}
\label{sec:prop1}

This section proves Proposition~1 of KL: how a flight-to-safety shock
$\dev{\cy}_1>0$ propagates to expected real rates, inflation, the relative price
of capital, and global employment. The proof specializes the
engine of Section~\ref{sec:engine} to identical portfolios and complete home
bias, which forces the wealth share to be inert, so that the only moving parts
are prices and production. We treat the two wage regimes in turn.

\subsection{Statement}

\begin{proposition}[Effects of a safety shock on prices and production]
\label{prop:prices}
Consider the simplified economy of Definition~\ref{def:simplified} with identical
portfolios across countries, no Home government safe debt
($b^g_{H,s}=0$), and complete home bias $\hbias\to\zs/(1+\zs)$, so that
$\dev{\wsh}_1=0$ and $\Eone\dev{\wsh}_2=0$. Consider a flight-to-safety shock
$\dev{\cy}_1>0$ in period~$1$, with $\widetilde{\widetilde{\dev k}}_1=0$ on
impact.
\begin{enumerate}[label=(\alph*),leftmargin=1.8em,itemsep=3pt]
  \item \textbf{Flexible wages} ($\frisch\to0$). The shock is absorbed entirely
  by lower expected real rates and Home deflation, with no effect on the real
  exchange rate or production:
  \begin{equation}
  \Eone\rr_2 = -\dev{\cy}_1, \qquad
  \D\dev{P}_1 = -\frac{1}{\taylor}\dev{\cy}_1, \qquad
  \dev{\rer}_1 = \dev{\ell}_1 = \fr{\dev{\ell}}_1 = 0. \label{eq:prop1-flex}
  \end{equation}
  \item \textbf{Wages set one period in advance.} The shock causes a global
  recession and the responses are \emph{strictly} attenuated relative to the
  flexible-wage benchmark:
  \begin{equation}
  0 > \Eone\rr_2 > -\dev{\cy}_1, \qquad
  0 > \D\dev{P}_1 > -\frac{1}{\taylor}\dev{\cy}_1, \qquad
  \dev{\rer}_1 \propto \dev{\cy}_1 \;(>0), \label{eq:prop1-sticky}
  \end{equation}
  and global employment falls, \emph{disproportionately at Home}:
  \begin{equation}
  \frac{1}{1+\zs}\dev{\ell}_1 + \frac{\zs}{1+\zs}\fr{\dev{\ell}}_1 \;\propto\; -\dev{\cy}_1, \qquad
  \dev{\ell}_1 - \fr{\dev{\ell}}_1 \;\propto\; -\dev{\cy}_1. \label{eq:prop1-labor}
  \end{equation}
\end{enumerate}
\end{proposition}

\noindent (We transcribe the proposition to match KL's main-paper statement; the
proportionality constants in \eqref{eq:prop1-sticky}--\eqref{eq:prop1-labor} are
made explicit in the proof.)

\subsection{Proof}

\paragraph{Specialization of the engine.} Following KL's app.~B.3.1, impose
identical portfolios --- so that the steady-state positions satisfy
$\qk\,k=\sav$, $b_F=0$, $b_H=0$ --- and no Home government safe debt,
$b^g_{H,s}=0$. The realized wealth-share equation~\eqref{eq:wsh-realized} then
has \emph{every} coefficient equal to zero:
\[
\dev{\wsh}_1 = \underbrace{\Bigl(\tfrac{\qk k}{\sav}-1\Bigr)}_{=0}\bigl(\rk_1-\rr_1\bigr)
+\underbrace{\tfrac{b_F}{\sav}}_{=0}\bigl(\fr{r}_1-\dev{\rer}_1-\rr_1\bigr)
-\disc\,\tfrac{\zs}{1+\zs}\underbrace{\tfrac{b^g_{H,s}}{\sav}}_{=0}\dev{\cy}_1
= 0.
\]
Moreover, with complete home bias $\hbias\to\zs/(1+\zs)$, the net Home
expenditure weight in the wealth-share law~\eqref{eq:Ewsh-solved} vanishes,
$\frac{\zs}{1+\zs}-\hbias\to0$, and the relative-labor channel multiplier
collapses; combined with $\dev{\wsh}_1=0$ this gives
\begin{equation}
\Eone\dev{\wsh}_2 = 0. \label{eq:prop1-wsh0}
\end{equation}
\gapflag{KL state ``since $\hbias\to\zs/(1+\zs)$, we have that $\Eone\dev{\wsh}_2=0$''
(their app.~B.3.1). Two facts combine. First, $\dev{\wsh}_1=0$ from the
identical-portfolio specialization above. Second, in \eqref{eq:Ewsh-solved} the
surviving relative-labor term carries the factor
$\bigl(\frac{\zs}{1+\zs}-\hbias\bigr)$ \emph{through} $\Eone\dev{\tot}_2$ and the
$\atil$ feedback: with complete home bias $\hbias=\frac{\zs}{1+\zs}$ one checks
$\hbias\frac{1+\zs}{\zs}=1$, so $1-(\hbias\frac{1+\zs}{\zs})^2=0$, killing the
relative-labor coefficient in \eqref{eq:Ewsh-solved}. Hence $\Eone\dev{\wsh}_2
=\dev{\wsh}_1=0$. We verified $\hbias\frac{1+\zs}{\zs}=1$ at complete home bias
and that this zeroes the $1-(\hbias\frac{1+\zs}{\zs})^2$ factor; KL suppress this
one-line check.}
With \eqref{eq:prop1-wsh0}, the wealth-share terms drop out of every reduced form
in Section~\ref{sec:engine-summary}. In particular, using
$\hbias\frac{1+\zs}{\zs}=1$ (complete home bias) the composite denominator
\eqref{eq:Lambda-def} simplifies to
\begin{equation}
\Lambda = \frac{\cshare}{1-\cshare}+\tel+(1-\tel)\cdot1
=\frac{\cshare}{1-\cshare}+1 = \frac{1}{1-\cshare}, \label{eq:Lambda-cmplt}
\end{equation}
so the relative-labor coefficient becomes
$\hbias/\Lambda=\frac{\zs}{1+\zs}(1-\cshare)$. We record this for the sticky-wage
case below.

\subparagraph{(a) Flexible wages.} With flexible wages and infinite Frisch
elasticity $\frisch\to0$, the union's wage-setting condition gives
\eqref{eq:flex-labor}:
\[
\dev{\ell}_1 = \fr{\dev{\ell}}_1 = 0.
\]
Substitute $\dev{\ell}_1=\fr{\dev{\ell}}_1=0$, $\Eone\dev{\wsh}_2=0$, and
$\widetilde{\widetilde{\dev k}}_1=0$ into the engine. The expected capital
return~\eqref{eq:Er2k-solved} becomes
\[
\Eone\rk_2 = -(1-\cshare)\underbrace{\Bigl(\tfrac{1}{1+\zs}\dev{\ell}_1+\tfrac{\zs}{1+\zs}\fr{\dev{\ell}}_1\Bigr)}_{=0}
+\frac{\hbias}{\Lambda}\underbrace{\bigl(\fr{\dev{\ell}}_1-\dev{\ell}_1\bigr)}_{=0} = 0.
\]
The capital-Euler relation~\eqref{eq:euler-r} then gives the expected real bond
return,
\begin{equation}
\Eone\rr_2 = \Eone\rk_2 - \dev{\cy}_1 = -\dev{\cy}_1, \label{eq:prop1a-Er2}
\end{equation}
which is the first claim in~\eqref{eq:prop1-flex}. The Home Taylor
rule~\eqref{eq:dP1} gives the inflation deviation,
\begin{equation}
\D\dev{P}_1 = \frac{1}{\taylor}\bigl(\Eone\rk_2-\dev{\cy}_1\bigr)
=\frac{1}{\taylor}\bigl(0-\dev{\cy}_1\bigr) = -\frac{1}{\taylor}\dev{\cy}_1, \label{eq:prop1a-dP}
\end{equation}
the second claim. Finally the capital-price relation~\eqref{eq:qk1-raw} with
$\Eone\dev{\tot}_2=0$ (from \eqref{eq:Es2} and \eqref{eq:prop1-wsh0}),
$\widetilde{\widetilde{\dev k}}_1=0$, and $\Eone\rk_2=0$ gives
\[
\dev{\qk}_1 = -\hbias\underbrace{\Eone\dev{\tot}_2}_{=0}-(1-\cshare)\underbrace{\widetilde{\widetilde{\dev k}}_1}_{=0}-\underbrace{\Eone\rk_2}_{=0}=0,
\]
which with \eqref{eq:flex-labor} completes \eqref{eq:prop1-flex}. KL note these
claims ``follow immediately from the above results given
$\dev{\ell}_1=\fr{\dev{\ell}}_1=0$.''

\subparagraph{(b) Wages set one period in advance.} Now wages are predetermined.
Substituting the inflation reduced forms~\eqref{eq:dP1}--\eqref{eq:dPstar1} (with
$\Eone\dev{\wsh}_2=0$) into the firm-labor-demand
system~\eqref{eq:firm-H}--\eqref{eq:firm-F}, and using complete home bias and
$\widetilde{\widetilde{\dev k}}_1=0$, KL reduce the $2\times2$ system to a linear
map from $(\dev{\cy}_1,\widetilde{\widetilde{\dev k}}_1)$ to
$(\dev{\ell}_1,\fr{\dev{\ell}}_1)$. The solution is (their app.~B.3.1)
\begin{align}
\dev{\ell}_1 &= -\frac{1}{\cshare+\frac{1}{\taylor}(1-\cshare)}\,\frac{1}{\taylor}\,\dev{\cy}_1
\;-\;\frac{1}{\cshare+\frac{1}{\taylor}(1-\cshare)}\,(1-\cshare)\,\widetilde{\widetilde{\dev k}}_1, \label{eq:prop1b-l}\\
\fr{\dev{\ell}}_1 &= -\frac{1}{\cshare+\frac{1}{\taylor}(1-\cshare)}\,(1-\cshare)\,\widetilde{\widetilde{\dev k}}_1. \label{eq:prop1b-lstar}
\end{align}
\gapflag{KL display \eqref{eq:prop1b-l}--\eqref{eq:prop1b-lstar} as the solution
of \eqref{eq:firm-H}--\eqref{eq:firm-F} but do not show the elimination. The
reconstruction: at complete home bias $\hbias\frac{1+\zs}{\zs}=1$, so the
terms-of-trade coefficients in \eqref{eq:firm-H}--\eqref{eq:firm-F} simplify
(the $\frac{\zs}{1+\zs}-\hbias$ and $\frac{1}{1+\zs}-\frac{\hbias}{\zs}$ pieces
each $\to0$, leaving only the $\frac{\cshare}{1-\cshare}$-weighted terms-of-trade
terms). Replacing $\D\dev P_1,\D\fr{\dev P}_1$ on the right via
\eqref{eq:dP1}--\eqref{eq:dPstar1} with $\Eone\rk_2$ from~\eqref{eq:Er2k-solved}
(and $\Lambda=1/(1-\cshare)$, so $\frac{\hbias}{\Lambda}=\frac{\zs}{1+\zs}(1-\cshare)$),
the system becomes, after collecting the global-employment combination
$L\equiv\frac{1}{1+\zs}\dev{\ell}_1+\frac{\zs}{1+\zs}\fr{\dev{\ell}}_1$ and the
relative-employment gap $G\equiv\dev{\ell}_1-\fr{\dev{\ell}}_1$:
\[
\cshare L + \tfrac{1}{\taylor}(1-\cshare)L = -\tfrac{1}{\taylor}\dev{\cy}_1
-(1-\cshare)\widetilde{\widetilde{\dev k}}_1,
\qquad
\bigl[\cshare+\tfrac{1}{\taylor}(1-\cshare)\bigr]\,(\text{relative part}) = -\tfrac{1}{\taylor}\dev{\cy}_1,
\]
where the first equation is the population-weighted (Home+Foreign) sum of
\eqref{eq:firm-H}--\eqref{eq:firm-F} and the second is the Home$-$Foreign
difference. The convenience shock $\dev{\cy}_1$ enters only the Home Taylor
rule~\eqref{eq:dP1} (Foreign's \eqref{eq:dPstar1} has no $\dev{\cy}_1$), so it
loads asymmetrically: it appears in the global sum and the difference but
\emph{not} in $\fr{\dev{\ell}}_1$ alone. Solving the two scalar equations for
$\dev{\ell}_1$ and $\fr{\dev{\ell}}_1$ gives
\eqref{eq:prop1b-l}--\eqref{eq:prop1b-lstar}. We checked: (i) $\fr{\dev{\ell}}_1$
has no $\dev{\cy}_1$ term --- the safety shock leaves Foreign employment
unaffected at complete home bias, consistent with KL's remark that
``$\hbias\to\frac{\zs}{1+\zs}$ implies $\fr{\dev{\ell}}_1$ is unaffected by a
safety shock''; (ii) the common denominator
$\cshare+\frac{1}{\taylor}(1-\cshare)$ matches KL's
$\alpha+\frac{1}{\phi}(1-\alpha)$. The forms match KL exactly with
$\frac{1}{\phi}=\frac{1}{\taylor}$.}

Set $\widetilde{\widetilde{\dev k}}_1=0$ on impact. Then
\eqref{eq:prop1b-l}--\eqref{eq:prop1b-lstar} reduce to
\begin{equation}
\dev{\ell}_1 = -\frac{1}{\taylor}\,\frac{1}{\cshare+\frac{1}{\taylor}(1-\cshare)}\,\dev{\cy}_1 < 0,
\qquad \fr{\dev{\ell}}_1 = 0. \label{eq:prop1b-impact}
\end{equation}
Because $\cshare\in(0,1)$ and $\taylor>1$, the denominator
$\cshare+\frac{1}{\taylor}(1-\cshare)>0$, so $\dev{\ell}_1<0$ for $\dev{\cy}_1>0$:
\emph{Home employment falls}, while Foreign employment is unchanged on impact.
We now read off the remaining claims.

\emph{Global employment falls disproportionately at Home.} From
\eqref{eq:prop1b-impact},
\begin{align}
\frac{1}{1+\zs}\dev{\ell}_1+\frac{\zs}{1+\zs}\fr{\dev{\ell}}_1
&= \frac{1}{1+\zs}\,\dev{\ell}_1
= -\frac{1}{1+\zs}\frac{1}{\taylor}\frac{1}{\cshare+\frac{1}{\taylor}(1-\cshare)}\dev{\cy}_1 \;\propto\; -\dev{\cy}_1, \\
\dev{\ell}_1-\fr{\dev{\ell}}_1 &= \dev{\ell}_1
= -\frac{1}{\taylor}\frac{1}{\cshare+\frac{1}{\taylor}(1-\cshare)}\dev{\cy}_1 \;\propto\; -\dev{\cy}_1,
\end{align}
establishing both proportionalities in~\eqref{eq:prop1-labor}: aggregate
employment falls (negative) and the Home$-$Foreign gap is negative (Home falls
by more), exactly the ``disproportionately at Home'' claim.

\emph{Expected capital return and expected real rate.} Substituting
\eqref{eq:prop1b-impact} into~\eqref{eq:Er2k-solved} with
$\Lambda=1/(1-\cshare)$:
\begin{align}
\Eone\rk_2 &= -(1-\cshare)\,\frac{1}{1+\zs}\dev{\ell}_1
+\frac{\zs}{1+\zs}(1-\cshare)\bigl(\fr{\dev{\ell}}_1-\dev{\ell}_1\bigr) \notag\\
&= -(1-\cshare)\frac{1}{1+\zs}\dev{\ell}_1 - \frac{\zs}{1+\zs}(1-\cshare)\dev{\ell}_1
= -(1-\cshare)\dev{\ell}_1 > 0, \label{eq:prop1b-Erk}
\end{align}
using $\fr{\dev{\ell}}_1=0$ and
$\frac{1}{1+\zs}+\frac{\zs}{1+\zs}=1$. Since $\dev{\ell}_1<0$, we have
$\Eone\rk_2>0$. Then the expected real bond return from~\eqref{eq:euler-r} is
\begin{equation}
\Eone\rr_2 = \Eone\rk_2 - \dev{\cy}_1 = -(1-\cshare)\dev{\ell}_1 - \dev{\cy}_1. \label{eq:prop1b-Er2}
\end{equation}
Substituting $\dev{\ell}_1$ from~\eqref{eq:prop1b-impact},
\[
\Eone\rr_2 = (1-\cshare)\frac{1}{\taylor}\frac{1}{\cshare+\frac{1}{\taylor}(1-\cshare)}\dev{\cy}_1 - \dev{\cy}_1
= -\dev{\cy}_1\left[1 - \frac{\frac{1}{\taylor}(1-\cshare)}{\cshare+\frac{1}{\taylor}(1-\cshare)}\right]
= -\dev{\cy}_1\,\frac{\cshare}{\cshare+\frac{1}{\taylor}(1-\cshare)}.
\]
The bracket $\frac{\cshare}{\cshare+\frac{1}{\taylor}(1-\cshare)}\in(0,1)$ for
$\cshare\in(0,1),\taylor>1$, so
\begin{equation}
0 > \Eone\rr_2 = -\frac{\cshare}{\cshare+\frac{1}{\taylor}(1-\cshare)}\dev{\cy}_1 > -\dev{\cy}_1, \label{eq:prop1b-Er2-ineq}
\end{equation}
which is the first strict-inequality chain in~\eqref{eq:prop1-sticky}: the
expected real rate falls, but by \emph{less} than the convenience yield ---
strictly between $0$ and $-\dev{\cy}_1$ --- because part of the safety shock is
now absorbed by the recession (the positive $\Eone\rk_2$) rather than by the real
rate alone.

\emph{Inflation.} From the Home Taylor rule~\eqref{eq:dP1},
\begin{equation}
\D\dev{P}_1 = \frac{1}{\taylor}\bigl(\Eone\rk_2-\dev{\cy}_1\bigr) = \frac{1}{\taylor}\Eone\rr_2
= -\frac{1}{\taylor}\frac{\cshare}{\cshare+\frac{1}{\taylor}(1-\cshare)}\dev{\cy}_1, \label{eq:prop1b-dP}
\end{equation}
so dividing the inequality~\eqref{eq:prop1b-Er2-ineq} by $\taylor>0$,
\[
0 > \D\dev{P}_1 > -\frac{1}{\taylor}\dev{\cy}_1,
\]
the second chain in~\eqref{eq:prop1-sticky}: Home still deflates, but less than in
the flexible benchmark~\eqref{eq:prop1a-dP}.

\emph{Real exchange rate.} At complete home bias $\hbias(1+\zs)/\zs=1$, so the
real-exchange-rate relation~\eqref{eq:q1-s1} reads $\dev{\rer}_1=\dev{\tot}_1$.
From the solved terms of trade~\eqref{eq:s1-solved} with $\Eone\dev{\wsh}_2=0$,
\begin{equation}
\dev{\rer}_1 = \dev{\tot}_1 = \frac{1}{\Lambda}\bigl(\fr{\dev{\ell}}_1-\dev{\ell}_1\bigr)
= -(1-\cshare)\dev{\ell}_1
= (1-\cshare)\frac{1}{\taylor}\frac{1}{\cshare+\frac{1}{\taylor}(1-\cshare)}\dev{\cy}_1 \;>\;0, \label{eq:prop1b-q}
\end{equation}
using $\Lambda=1/(1-\cshare)$ at complete home bias and $\fr{\dev{\ell}}_1=0$,
$\dev{\ell}_1<0$. Since the convention is that a rise in $\rer$ is a \emph{Home
appreciation}, a positive safety shock \emph{appreciates} the Home real exchange
rate by a positive multiple of $\dev{\cy}_1$, establishing the third claim
in~\eqref{eq:prop1-sticky} (in the flexible case $\dev{\ell}_1=0$ so
$\dev{\rer}_1=0$, the third claim in~\eqref{eq:prop1-flex}). \qed
(\,Proposition~\ref{prop:prices}\,)

\begin{remark}[The relative price of capital, an engine by-product]
KL's Proposition~1 makes no claim about the real price of installed capital
$\qk$, but the engine pins it down and it is useful downstream (Propositions~2--3).
From the capital-price relation~\eqref{eq:qk1-raw} with $\Eone\dev{\tot}_2=0$ and
$\widetilde{\widetilde{\dev k}}_1=0$,
\begin{equation}
\dev{\qk}_1 = -\Eone\rk_2 = (1-\cshare)\dev{\ell}_1
= -(1-\cshare)\frac{1}{\taylor}\frac{1}{\cshare+\frac{1}{\taylor}(1-\cshare)}\dev{\cy}_1 \;<\;0. \label{eq:prop1b-qk}
\end{equation}
The price of installed capital \emph{falls} in a flight to safety: the expected
capital return $\Eone\rk_2$ rises (the recession raises the marginal product of
capital), and the no-arbitrage condition $\dev{\qk}_1=-\Eone\rk_2$ revalues
capital down. Note $\dev{\qk}_1=-\dev{\rer}_1$ here: the price of capital and the
real exchange rate move by equal and opposite amounts on impact.
\end{remark}

\subsection{Discussion}

\paragraph{The Euler-equation wedge.} The entire result is driven by the
convenience term in the Home Euler equation. The safe bond enters the
intertemporal condition through the factor $(1+r_{t+1})/(1-\cy_t)$: holders accept
a lower pecuniary return because the bond yields nonpecuniary (safety/liquidity)
services. To first order this is the wedge~\eqref{eq:euler-r},
$\Eone\rk_2-\Eone\rr_2=\dev{\cy}_1$: a flight to safety raises the required
\emph{spread} of capital over safe bonds by exactly $\dev{\cy}_1$.

\paragraph{Flexible wages: one-for-one absorption.} With flexible wages and
prices, output cannot move ($\dev{\ell}_1=\fr{\dev{\ell}}_1=0$, so the marginal
product of capital and hence $\Eone\rk_2$ are unchanged). The required spread is
then met \emph{entirely} by a fall in the expected real rate,
$\Eone\rr_2=-\dev{\cy}_1$. Under the Taylor rule, a lower expected real rate
requires Home deflation, $\D\dev{P}_1=-\frac{1}{\taylor}\dev{\cy}_1$: the central
bank delivers exactly enough deflation that the Fisher relation produces the
needed real-rate decline. There is no recession and no revaluation of capital.

\paragraph{Sticky wages: a deflationary recession, worse at Home.} When nominal
wages are predetermined, the deflationary pressure from the flight to safety
raises the \emph{product wage} $W/P$ (since $P$ falls but $W$ is fixed), which
depresses labor demand and output. The recession raises the expected marginal
product of capital, so $\Eone\rk_2>0$ absorbs part of the required spread; the
expected real rate therefore falls by \emph{less} than $\dev{\cy}_1$
(eq.~\eqref{eq:prop1b-Er2-ineq}), and deflation is correspondingly milder
(eq.~\eqref{eq:prop1b-dP}). The recession is concentrated at Home because the
safety shock enters only the \emph{Home} Taylor rule: the convenience yield is on
\emph{dollar} bonds, so it is Home monetary policy that must accommodate it.
Foreign employment is unchanged on impact ($\fr{\dev{\ell}}_1=0$), so Home's
terms of trade and (via~\eqref{eq:q1-s1}) the real exchange rate \emph{appreciate}
($\dev{\ell}_1<\fr{\dev{\ell}}_1=0\Rightarrow\dev{\tot}_1=\frac{1}{\Lambda}(\fr{\dev{\ell}}_1-\dev{\ell}_1)>0$):
the dollar appreciates in the flight to safety.

\paragraph{Load-bearing assumptions.} (i) \emph{Nominal rigidity} is essential
for the recession: with flexible wages the real allocation is untouched and only
prices move. (ii) The \emph{Taylor rule that does not cut nominal rates
one-for-one with $\cy$} (an inflation coefficient $\taylor$, not a peg on the
convenience-adjusted rate) is what forces deflation to do the work; a central
bank that fully accommodated the safety shock would short-circuit the result.
(iii) The \emph{trade elasticity} $\tel$ governs the size of the terms-of-trade /
real-exchange-rate response through $\Lambda$; here, at complete home bias,
$\Lambda=1/(1-\cshare)$ and the $\tel$-dependence drops from the labor solution,
but it returns away from complete home bias. (iv) \emph{Complete home bias}
$\hbias\to\zs/(1+\zs)$ is what delivers the clean $\dev{\wsh}_1=\Eone\dev{\wsh}_2=0$,
isolating the price/production channel from wealth-share redistribution (which is
the subject of Propositions~\ref{prop:returns}--\ref{prop:wealth}).

\paragraph{Connection to the target facts.} Proposition~\ref{prop:prices} is the
model's account of the empirical regularity that a global flight to safety
coincides with (a) \emph{dollar appreciation} --- here Home's terms of trade and
real exchange rate appreciate because the dollar convenience shock forces Home
deflation and a Home-concentrated contraction --- and (b) a \emph{global
recession that hits the safe-asset issuer's economy hardest}. In the big\_cy
reading, the same mechanism with a \emph{large absolute} convenience yield
$\dev{\cy}_1$ amplifies both the appreciation and the asymmetric contraction,
which is the channel we will quantify downstream.
